%% IEEE IoTJ submission template
%% Title: Dual-Constrained DRL for Joint UAV Trajectory Design and
%%        NOMA-Based Task Offloading in MEC Networks

\documentclass[journal]{IEEEtran}

\usepackage{amsmath,amssymb,amsfonts,amsthm}
\newtheorem{lemma}{Lemma}
\newtheorem{theorem}{Theorem}
\newtheorem{corollary}{Corollary}
\newtheorem{remark}{Remark}
\usepackage{graphicx}
\usepackage{algorithmic}
\usepackage{algorithm}
\usepackage{booktabs}
\usepackage{array}
\usepackage{multirow}
\usepackage{cite}
\usepackage{url}
\usepackage{xcolor}
\usepackage{bm}

%% ── Shortcuts ──────────────────────────────────────────────────────────────
\newcommand{\uav}{\text{UAV}}
\newcommand{\mec}{\text{MEC}}
\newcommand{\noma}{\text{NOMA}}
\newcommand{\RR}{\mathbb{R}}
\newcommand{\EE}{\mathbb{E}}
\newcommand{\calC}{\mathcal{C}}
\newcommand{\calT}{\mathcal{T}}
\newcommand{\ld}{\lambda_d}
\newcommand{\todo}[1]{\textcolor{red}{[TODO: #1]}}

\begin{document}

\title{Dual-Constrained Deep Reinforcement Learning for Joint UAV
Trajectory Design and NOMA-Based Task Offloading in MEC Networks}

\author{
  \IEEEauthorblockN{Mingxi Guan}
  \IEEEauthorblockA{%
  School of Computer Science and Engineering\\
  (Affiliation to be finalized before submission)}
}

\maketitle

%% ═══════════════════════════════════════════════════════════════════════════
\begin{abstract}
Unmanned aerial vehicle (UAV)-assisted mobile edge computing (MEC)
with non-orthogonal multiple access (NOMA) offers a promising
paradigm for low-latency IoT task offloading in dynamic environments.
However, jointly optimizing UAV trajectory, NOMA power allocation,
and offloading decisions under hard deadline constraints is
NP-hard and difficult to solve with conventional methods.
In this paper, we formulate the problem as a constrained Markov
decision process (CMDP) and propose \textbf{DC-PPO}
(Dual-Constrained Proximal Policy Optimization), which incorporates
a Lagrangian dual variable with a physically-derived upper bound
$\lambda_{\max}$ to enforce delay constraints without destabilizing
the on-policy training dynamics.
The key insight of DC-PPO is a reward-cost ratio upper bound
$\lambda_{\text{crit}} = \bar{r}/\bar{c}$, which ensures the
augmented reward remains positive throughout training.
Extensive experiments with $N{=}20$ users demonstrate that
DC-PPO achieves a deadline satisfaction rate of \textbf{82.5\%},
outperforming unconstrained PPO by \textbf{7.3 percentage points}
(p\,=\,0.031, Cohen's $d$\,=\,2.69) and reducing training
variance by 47\% compared to PPO-Lagrangian.
Scalability experiments confirm that the performance advantage
grows monotonically with user density, reaching $+6.8$\,pp
at $N{=}25$.
An unexpected finding is the \emph{NOMA interference saturation
effect}: myopic full-offloading achieves only 15.4\% deadline
satisfaction---far below random (63.0\%)---because simultaneous
maximum-power transmission collapses the NOMA SINR for all users.
\end{abstract}

\begin{IEEEkeywords}
UAV-MEC, NOMA, task offloading, constrained Markov decision process,
proximal policy optimization, Lagrangian relaxation, deadline constraint.
\end{IEEEkeywords}

%% ═══════════════════════════════════════════════════════════════════════════
\section{Introduction}\label{sec:intro}

\IEEEPARstart{T}{he} proliferation of compute-intensive IoT
applications---ranging from real-time video analytics to
autonomous vehicle coordination---is straining the processing
capacity of resource-constrained edge devices.
Mobile Edge Computing (MEC) addresses this by deploying
computation servers at the network edge, enabling low-latency
task offloading \cite{mao2017mec_survey,shi2016mec_vision}.
However, deploying MEC infrastructure at fixed base stations
is impractical in dynamic or disaster-affected environments
where ground connectivity is disrupted.
Unmanned aerial vehicles (UAVs) equipped with MEC servers offer
a compelling complement: their mobility allows on-demand
repositioning to serve user clusters, and line-of-sight (LoS)
channels to ground devices yield significantly higher
transmission rates than terrestrial links
\cite{zeng2019uav_survey,mozaffari2019uav_tutorial}.

UAV-assisted MEC faces a fundamental throughput bottleneck
when serving multiple users simultaneously:
traditional orthogonal multiple access (OMA) schemes allocate
disjoint frequency bands, limiting spectral efficiency as the
number of users grows.
Non-orthogonal multiple access (NOMA), which superposes
multiple users on the same channel and relies on successive
interference cancellation (SIC) at the receiver, can serve
$N$ users simultaneously within the same bandwidth
\cite{ding2017noma,liu2017noma_proc,ding2017noma_survey},
making it a natural fit for the dense IoT offloading regime.

\noindent\textbf{Challenges and Motivation.}

Jointly optimizing \emph{UAV trajectory}, \emph{NOMA power allocation},
and \emph{computation offloading ratio} presents three
inter-coupled challenges:

\textbf{C1 -- High-dimensional continuous action space.}
Even for $N{=}20$ users, the joint action space has 42 continuous
dimensions (UAV displacement + per-user offloading ratio +
per-user NOMA power), ruling out discrete-action RL methods
such as DQN and making tabular or convex approaches intractable.

\textbf{C2 -- Hard deadline constraints.}
IoT tasks carry strict latency deadlines (e.g., $t_{\max}{=}1.0$\,s
for compressed image analytics), yet the delay depends
non-linearly on the NOMA interference structure, which varies
with UAV position and user power levels.
Existing unconstrained RL methods cannot reliably enforce such
constraints during training.

\textbf{C3 -- Lagrangian instability.}
The Lagrangian relaxation approach \cite{tessler2019rcpo} converts
the constrained optimization into a penalized objective, but the
dual multiplier $\ld$ must be carefully bounded.
Without a principled upper bound, $\ld$ can grow without limit,
causing the augmented reward to become negative and collapsing
the policy gradient---a failure mode we term
\emph{Lagrangian saturation}.

\noindent\textbf{Related Work.}

\textbf{UAV trajectory optimization.}
Conventional approaches formulate UAV-MEC as a mixed-integer
non-linear program and apply successive convex approximation
(SCA) or block coordinate descent
\cite{zeng2019uav_survey,zhan2020uav_mec,xie2020uav_mec_joint,
cao2021uav_mec_offload}.
These methods yield locally optimal solutions but require
accurate channel state information and cannot handle stochastic
task arrivals without recomputation each slot.

\textbf{Deep RL for MEC offloading.}
DRL methods learn offloading policies directly from environment
interactions, avoiding the need for explicit channel modeling
\cite{luong2019drl_comms,mnih2015dqn}.
DDPG and TD3 handle continuous action spaces but require
large replay buffers and struggle in high-dimensional settings
under limited training budgets \cite{fujimoto2018td3}.
On-policy methods such as PPO offer more stable training but
lack a principled mechanism for constraint enforcement.
DC-PPO uses identical hyperparameters and network architecture
to its PPO baseline (128 hidden, Tanh, orthogonal init,
obs normalization, entropy coeff.\ 0.01, clip $\epsilon{=}0.2$);
the sole difference is the Lagrangian mechanism.

\textbf{Constrained RL.}
Reward Constrained Policy Optimization (RCPO) \cite{tessler2019rcpo}
and Constrained Policy Optimization (CPO) \cite{achiam2017cpo}
provide theoretical frameworks for CMDP-based learning.
CPO relies on a computationally expensive trust-region projection
that scales poorly to 42-dimensional action spaces.
RCPO's Lagrangian approach is more practical, but it does not
specify how to bound the multiplier, leading to the saturation
problem identified in this work.
More recent methods include FOCOPS \cite{zhang2020focops}, which
constrains the first-order policy update, and CUP
\cite{yang2022cup}, which uses a constrained update projection.
However, these methods require computing the constraint cost
gradient, which adds significant overhead in high-dimensional
continuous action spaces.
PID-Lagrangian methods \cite{stooke2020pid_lagrangian} apply
proportional-integral-derivative control to the multiplier update
to reduce oscillation, but still lack a principled upper bound.
To our knowledge, no prior work has derived a
\emph{physically-grounded} upper bound for the Lagrangian
multiplier in the UAV-MEC context.

\textbf{UAV-NOMA-MEC.}
Recent works have combined UAV trajectory design with NOMA
offloading.
Wang~\emph{et al.} \cite{wang2023uav_mec_drl} applied DRL
to multi-UAV MEC task scheduling, and
Zhang~\emph{et al.} \cite{zhang2023uav_noma_iot} studied
energy-efficient UAV-NOMA for IoT with imperfect CSI.
Chen~\emph{et al.} \cite{chen2024uav_mec_multi} addressed
task-dependent delay in multi-UAV MEC scheduling.
Hu \emph{et al.} \cite{hu2023uav_noma_mec} jointly optimized
UAV trajectory and offloading ratio for NOMA-MEC networks via
SCA, demonstrating throughput gains over OMA baselines.
Li \emph{et al.} \cite{li2022noma_uav_drl} applied DRL to
NOMA-assisted UAV-MEC resource allocation but used an
unconstrained reward without deadline enforcement.
In contrast, DC-PPO explicitly handles deadline constraints
via a physically-grounded Lagrangian bound, and provides
statistical validation across 5 seeds with a 42-dimensional
joint action space---a setting not addressed by prior work.

\noindent\textbf{Contributions.}

This paper proposes \textbf{DC-PPO} (Dual-Constrained PPO),
which addresses the above challenges through three contributions:

\begin{itemize}
  \item \textbf{C1 [CMDP Formulation]:}
    We provide an integrated CMDP formulation that jointly
    couples UAV trajectory, per-user NOMA power allocation, and
    partial offloading ratio in a single 42-dimensional
    continuous action space, with an episode-level delay
    constraint explicitly capturing the NOMA interference
    structure under SIC decoding.

  \item \textbf{C2 [$\lambda_{\max}$ Selection Criterion]:}
    We derive a \emph{reward-cost ratio upper bound}
    $\lambda_{\text{crit}} = \bar{r}/\bar{c}$ that guarantees
    the Lagrangian-augmented reward remains positive throughout
    training (Lemma~\ref{lem:rpc}).
    Setting $\lambda_{\max}$ to a physically-derived fraction of
    $\lambda_{\text{crit}}$ prevents \emph{Lagrangian saturation}
    (the collapse of the policy gradient when $\ld$ exceeds
    $\lambda_{\text{crit}}$), reducing training variance by 47\%
    compared to unconstrained PPO-Lagrangian.

  \item \textbf{C3 [Empirical Validation]:}
    Extensive experiments with $N{=}20$ users and 5 independent
    seeds demonstrate that DC-PPO achieves 82.5\% deadline
    satisfaction, outperforming unconstrained PPO by
    7.3\,percentage points ($p{=}0.031$, Cohen's $d{=}2.69$).
    Scalability experiments confirm that the DC-PPO advantage
    grows monotonically with user density, reaching $+6.8$\,pp
    at $N{=}25$.
\end{itemize}

The remainder of this paper is organized as follows.
Section~\ref{sec:model} presents the system model and CMDP
formulation.
Section~\ref{sec:algorithm} describes DC-PPO, including the
$\lambda_{\max}$ derivation and algorithm pseudocode.
Section~\ref{sec:setup} details the experimental setup.
Section~\ref{sec:results} presents and analyzes the results.
Section~\ref{sec:conclusion} concludes the paper.

%% ═══════════════════════════════════════════════════════════════════════════
\section{System Model and Problem Formulation}\label{sec:model}

\subsection{Network Topology}

We consider a UAV-assisted MEC system operating over $T$ discrete
time slots $\mathcal{T} = \{1, \ldots, T\}$, each of duration
$\Delta t = 1$\,s.
A single UAV equipped with an MEC server serves $N$ ground IoT
users $\mathcal{N} = \{1, \ldots, N\}$ distributed uniformly at
random in a square area of side $L = 200$\,m.
The UAV flies at a fixed altitude $H = 100$\,m and communicates
with users via NOMA uplink on a bandwidth of $B = 5$\,MHz.
Let $\bm{q}_t = (x_t, y_t)^{\top} \in [0,L]^2$ denote the
UAV's horizontal position at slot $t$.

\subsection{UAV Mobility Model}

The UAV moves horizontally between consecutive slots.
The action $\bm{a}^{(\text{fly})}_t = (\delta x_t, \delta y_t)^{\top}$
specifies the desired displacement, subject to a maximum speed
constraint:
\begin{equation}
  \|\bm{a}^{(\text{fly})}_t\|_2 \le v_{\max} \cdot \Delta t,
  \label{eq:speed}
\end{equation}
with $v_{\max} = 20$\,m/s.
The UAV position is updated as:
\begin{equation}
  \bm{q}_{t+1} = \text{clip}\!\left(
    \bm{q}_t + \bm{a}^{(\text{fly})}_t,\; [0,L]^2
  \right).
  \label{eq:mobility}
\end{equation}
The propulsion energy consumed per slot is modeled as
\cite{zeng2019uav_survey}:
\begin{equation}
  E^{(\text{fly})}_t = \left(c_1 + c_2 v_t^2\right)\Delta t,
  \label{eq:fly_energy}
\end{equation}
where $v_t = \|\bm{a}^{(\text{fly})}_t\|_2 / \Delta t$ is the
flight speed, $c_1 = 80$\,W is the hovering power, and
$c_2 = 0.12$\,W$\cdot$s$^2$/m$^2$ is a velocity-dependent coefficient.

\subsection{Air-to-Ground NOMA Channel}

\subsubsection{Channel Model}
The air-to-ground channel from user $i$ to the UAV follows a
probabilistic LoS model \cite{zeng2019uav_survey}.
Let $d_i = \|\bm{q}_t - \bm{u}_i\|_2$ be the horizontal distance,
where $\bm{u}_i \in \mathbb{R}^2$ is user $i$'s fixed location.
The LoS probability is:
\begin{equation}
  P_{\text{LoS},i} = \frac{1}{1 + a \exp\!\left(-b(\theta_i - a)\right)},
  \label{eq:plos}
\end{equation}
where $\theta_i = \arctan(H / d_i)$ (in degrees) is the elevation
angle, and $a = 9.61$, $b = 0.16$ are environment-specific
parameters for urban areas \cite{zeng2019uav_survey}.
The average path loss (linear) is:
\begin{equation}
  \text{PL}_i = \text{PL}^{\text{fs}}_i \cdot
    \left(P_{\text{LoS},i} \cdot \eta_{\text{LoS}} +
          (1 - P_{\text{LoS},i}) \cdot \eta_{\text{NLoS}}\right),
  \label{eq:pl}
\end{equation}
where $\text{PL}^{\text{fs}}_i = (4\pi f_c \ell_i / c)^2$
is the free-space path loss with $\ell_i = \sqrt{d_i^2 + H^2}$,
$f_c = 2$\,GHz, and $\eta_{\text{LoS}} = 1$\,dB,
$\eta_{\text{NLoS}} = 20$\,dB are the additional attenuation
factors for LoS and NLoS links, respectively
(linear equivalents: $10^{0.1} \approx 1.26$ and $10^{2} = 100$).
The channel power gain is $h_i = 1/\text{PL}_i$.

\subsubsection{NOMA Uplink with SIC}

In each slot, active users transmit simultaneously on the shared
5\,MHz channel under NOMA.
Let $p_i \in [0, p_{\max}]$ be the transmit power of user $i$,
with $p_{\max} = 0.2$\,W.
The UAV receiver applies Successive Interference Cancellation (SIC).
Let $\pi(1), \pi(2), \ldots, \pi(N)$ denote the SIC decoding order
sorted by \emph{descending} received power, i.e.,
$p_{\pi(1)} h_{\pi(1)} \ge p_{\pi(2)} h_{\pi(2)} \ge \cdots$.
User $\pi(k)$ is decoded $k$-th, treating users
$\pi(k{+}1), \ldots, \pi(N)$ as interference (users already decoded
are perfectly cancelled).
The SINR of the $k$-th decoded user is:
\begin{equation}
  \text{SINR}_{\pi(k)} = \frac{p_{\pi(k)} h_{\pi(k)}}
    {\displaystyle\sum_{m=k+1}^{N} p_{\pi(m)} h_{\pi(m)} + \sigma^2},
  \label{eq:sinr}
\end{equation}
where $\sigma^2 = N_0 B = -110$\,dBm is the total noise power
over the $B{=}5$\,MHz bandwidth, with noise power spectral
density $N_0 = -174$\,dBm/Hz.
The achievable uplink rate is:
\begin{equation}
  R_i = B \log_2(1 + \text{SINR}_i) \quad \text{(bits/s)}.
  \label{eq:rate}
\end{equation}

\subsection{Task Model and Computation Offloading}

\subsubsection{Task Arrival}
In each slot $t$, user $i$ generates a computation task independently
with probability $p_a = 0.5$.
A task is parameterized by its data size $D_i \sim \mathcal{U}[0.4, 2.0]$\,Mbit
and computation density $\kappa_i \sim \mathcal{U}[500, 1000]$\,cycles/bit,
so the required computation is $C_i = D_i \kappa_i$ cycles.

\subsubsection{Partial Offloading}
For each active user $i$, the agent decides an offloading ratio
$\alpha_i \in [0,1]$ (fraction of task bits offloaded to UAV-MEC)
and an NOMA transmit power $p_i \in [0, p_{\max}]$.
The offloaded fraction $\alpha_i D_i$ is transmitted to the UAV
at rate $R_i$ and executed on the MEC server at an allocated
CPU frequency $f_i^{(\mec)}$.
The remaining fraction $(1-\alpha_i) D_i$ is processed locally at
CPU frequency $f_{\text{loc}} = 1$\,GHz.

The end-to-end task completion delay for user $i$ (denoted
$\tau_i$ to distinguish from the horizontal distance $d_i$) is:
\begin{equation}
  \tau_i = \max\!\left(
    \frac{(1-\alpha_i)C_i}{f_{\text{loc}}},\;
    \frac{\alpha_i D_i}{R_i} + \frac{\alpha_i C_i}{f_i^{(\mec)}}
  \right),
  \label{eq:delay}
\end{equation}
where the $\max$ reflects parallel execution of local and
offload portions (a standard assumption in partial offloading
models \cite{liu2021uav_mec_survey}).
We neglect MEC queuing delay, which is reasonable under the
equal-frequency partitioning assumption.
We adopt \emph{equal frequency partitioning} as a tractable
approximation that focuses the optimization on trajectory and
offloading decisions:
$f_i^{(\mec)} = f_{\uav} / N_{\text{active},t}$,
where $f_{\uav} = 8$\,GHz and $N_{\text{active},t}$ is the number
of active users in slot $t$.
We note that variable frequency allocation could be included as
an additional action dimension at the cost of a larger action space;
we leave this extension for future work.

The user-side energy consumption is:
\begin{equation}
  e_i = \kappa_{\text{e}} f_{\text{loc}}^2 \cdot (1-\alpha_i) C_i
      + p_i \cdot \frac{\alpha_i D_i}{R_i},
  \label{eq:energy}
\end{equation}
where the first term is the local computation energy
with effective capacitance coefficient
$\kappa_{\text{e}} = 10^{-27}$\,J\,$\cdot$\,s$^2$/cycle$^2$
(so that $\kappa_{\text{e}} f_{\text{loc}}^2 C$ has units of Joules),
and the second term is the transmission energy.
A task is deemed to meet its deadline if $\tau_i \le t_{\max} = 1.0$\,s.

\subsection{Problem Formulation as CMDP}

The joint optimization of UAV trajectory, NOMA power allocation,
and offloading ratio is formulated as a Constrained Markov Decision
Process (CMDP) \cite{altman1999cmdp}.

\subsubsection{State Space}
The state at slot $t$ is:
\begin{equation}
  \bm{s}_t = \left[\frac{\bm{q}_t}{L},\;
    \left\{d_i/L,\; \mathbb{1}[i\text{ active}],\;
           D_i/D_{\max},\; C_i/C_{\max},\;
           h_i'\right\}_{i=1}^N,\;
    \frac{T-t}{T}\right] \in \mathbb{R}^{2+5N+1},
  \label{eq:state}
\end{equation}
yielding a 103-dimensional state for $N = 20$,
where $d_i$ is the horizontal UAV-to-user distance,
$h_i' = \log_{10}(h_i + 10^{-12})$ is the log-normalized channel
gain, and $D_{\max}$, $C_{\max}$ are normalization constants.
User positions $\bm{u}_i$ are fixed within an episode and
re-sampled uniformly at each episode reset, so the agent must
generalize over different user distributions.

\subsubsection{Action Space}
The continuous action at each slot is:
\begin{equation}
  \bm{a}_t = \left[\delta x_t,\; \delta y_t,\;
    \{\alpha_{i,t}\}_{i=1}^N,\; \{p_{i,t}\}_{i=1}^N\right]
    \in \mathbb{R}^{2+2N},
  \label{eq:action}
\end{equation}
a 42-dimensional continuous space for $N = 20$.
All components are mapped to $[-1,1]$ via the neural network
output and linearly scaled to their physical ranges.

\subsubsection{Reward Function}
The slot-level reward balances task throughput against UAV
propulsion energy:
\begin{equation}
  r_t = \sum_{i \in \mathcal{A}_t} w_i
        - \beta E^{(\text{fly})}_t / E_{\max},
  \label{eq:reward}
\end{equation}
where $\mathcal{A}_t$ is the set of active users in slot $t$,
$w_i = 1$ rewards the agent for each task handled
(irrespective of deadline status, following the CMDP paradigm
in which the objective and constraint are separated),
$\beta = 0.01$ is a weighting coefficient that keeps the
flight energy penalty small relative to task throughput, and
$E_{\max} = (c_1 + c_2 v_{\max}^2)\Delta t = 128$\,J is the
maximum single-slot flight energy (at full speed), used to
normalize the energy term to $[0,1]$.
Deadline satisfaction is enforced exclusively through the
Lagrangian constraint in \eqref{eq:cstr}.

\subsubsection{Constraint}

We impose a soft delay constraint via a Lagrangian penalty.
Let $\bar{\tau}_t \triangleq |\mathcal{A}_t|^{-1}
\sum_{i \in \mathcal{A}_t} \tau_{i,t}$ denote the per-slot
average task delay.
The episode-level average delay constraint is:
\begin{equation}
  \mathbb{E}_\pi\!\left[\frac{1}{T}\sum_{t=1}^T \bar{\tau}_t\right]
    \le \tau^*,
  \label{eq:constraint}
\end{equation}
where $\tau^* = 0.5$\,s/task is the delay budget.
The hard deadline is $t_{\max} = 1.0$\,s; the soft budget
$\tau^* = 0.5$\,s provides an anticipatory signal that guides
the policy to avoid deadline violations before they occur.

\subsubsection{CMDP Objective}
The optimization problem is:
\begin{align}
  \max_{\pi} \quad & \mathbb{E}_{\pi}\!\left[
    \sum_{t=1}^{T} \gamma^{t-1} r_t \right] \label{eq:obj} \\
  \text{s.t.} \quad & \mathbb{E}_{\pi}\!\left[
    \frac{1}{T}\sum_{t=1}^T \bar{\tau}_t \right] \le \tau^*,
    \quad \text{(average delay budget)} \label{eq:cstr} \\
  & \eqref{eq:speed},\, \eqref{eq:mobility},\,
    \alpha_{i,t} \in [0,1],\, p_{i,t} \in [0, p_{\max}],
    \notag
\end{align}
where $\pi$ denotes the policy and $\gamma = 0.99$ is the
discount factor.
The CMDP \eqref{eq:obj}--\eqref{eq:cstr} is computationally
intractable in general.

\begin{lemma}[NP-Hardness of Single-Slot Subproblem]
\label{lem:nphard}
The single-slot subproblem of \eqref{eq:obj} (fixed UAV
position, equal power $p_i {=} p_{\max}$, binary offloading
$\alpha_i {\in} \{0,1\}$) is NP-hard.
\end{lemma}
\begin{proof}
We reduce from the \emph{Maximum Feasible Subsystem} (MaxFS)
problem, which is NP-hard \cite{luo2008dynamic}.
With equal power $p_i {=} p_{\max}$ and binary $\alpha_i$,
the rate $R_i$ of user $i$ (if offloading) depends on the
SIC decoding order of the offloading subset
$\mathcal{S} = \{i : \alpha_i {=} 1\}$.
User $i \in \mathcal{S}$ meets its deadline iff:
$D_i / R_i(\mathcal{S}) + C_i / f_i^{(\mec)} \le t_{\max}$,
where $R_i(\mathcal{S})$ is a non-monotone function of
$\mathcal{S}$ through SINR coupling \eqref{eq:sinr}.
Finding the largest $\mathcal{S}$ such that all users in
$\mathcal{S}$ meet their deadlines is equivalent to
MaxFS over the SINR-rate-delay inequality system,
which is NP-hard for $N \ge 3$ \cite{luo2008dynamic}.
The continuous relaxations ($\alpha_i \in [0,1]$, variable
$p_i$) do not reduce the complexity class.
\end{proof}

\noindent Since the single-slot subproblem is NP-hard and the
full CMDP additionally optimizes over UAV trajectory (coupling
channel gains $h_i$ across slots), the multi-slot problem
\eqref{eq:obj}--\eqref{eq:cstr} is at least as hard,
motivating the deep reinforcement learning approach proposed
in Section~\ref{sec:algorithm}.

%% ═══════════════════════════════════════════════════════════════════════════
\section{Proposed DC-PPO Algorithm}\label{sec:algorithm}

\subsection{Lagrangian Relaxation of the CMDP}

We solve the CMDP \eqref{eq:obj}--\eqref{eq:cstr} via the
Lagrangian dual method \cite{altman1999cmdp,tessler2019rcpo}.
Introducing a non-negative dual variable (Lagrangian multiplier)
$\ld \ge 0$ for the delay constraint \eqref{eq:cstr}, the
primal problem is converted into an unconstrained penalized
objective:
\begin{equation}
  \mathcal{L}(\pi, \ld) \triangleq
  \mathbb{E}_\pi\!\left[\sum_{t=1}^T \gamma^{t-1} r_t\right]
  - \ld \cdot \mathbb{E}_\pi\!\left[
      \frac{1}{T}\sum_{t=1}^T \bar{\tau}_t - \tau^*
    \right].
  \label{eq:lagrangian}
\end{equation}
At each time step the agent observes the augmented reward:
\begin{equation}
  r'_t = r_t - \ld \cdot c_t, \quad c_t = \bar{\tau}_t,
  \label{eq:aug_reward}
\end{equation}
where $c_t = \bar{\tau}_t$ is the per-task average delay in slot $t$.
The multiplier $\ld$ is updated via projected gradient ascent
after each episode:
\begin{equation}
  \ld \leftarrow \Pi_{[0,\,\lambda_{\max}]}\!\left(
    \ld + \eta_\lambda \left(\bar{\tau}_{\text{ep}} - \tau^*\right)
  \right),
  \label{eq:lambda_update}
\end{equation}
where $\bar{\tau}_{\text{ep}}$ is the episode-averaged per-task delay,
$\eta_\lambda = 3\times10^{-2}$ is the dual learning rate, and
$\Pi_{[\cdot]}$ denotes projection onto $[0, \lambda_{\max}]$.

\subsection{Physical Upper Bound for $\lambda_{\max}$}

The choice of $\lambda_{\max}$ is the key design parameter of DC-PPO.
Setting $\lambda_{\max}$ too small renders the constraint inactive;
setting it too large causes the augmented reward $r'_t$ to become
negative, collapsing the policy gradient direction.

\begin{lemma}[Policy Gradient Collapse under Lagrangian Saturation]
\label{lem:rpc}
Let $\bar{r} \triangleq \mathbb{E}_\pi[r_t]$ and
$\bar{c} \triangleq \mathbb{E}_\pi[c_t]$ denote the steady-state
mean reward and cost under a feasible policy $\pi$.
Define the \emph{critical multiplier}
$\lambda_{\text{crit}} \triangleq \bar{r}/\bar{c}$.
If $\ld > \lambda_{\text{crit}}$, then:
\begin{enumerate}
  \item[(i)] $\mathbb{E}_\pi[r'_t] < 0$, i.e., every action yields
    negative expected augmented return;
  \item[(ii)] the policy gradient magnitude satisfies
    $\|\nabla_\theta J'(\theta)\| \le
      \frac{2 \gamma T}{(1-\gamma)^2}
      (\sigma_r + \ld \sigma_c)\,
      \|\nabla_\theta \log \pi_\theta\|_\infty$,
    where $\sigma_r^2 = \text{Var}[r_t]$ and
    $\sigma_c^2 = \text{Var}[c_t]$.
    As $\ld \to \infty$, the gradient direction is dominated
    by the cost term, and \emph{all actions are penalized
    uniformly regardless of task completion quality}.
\end{enumerate}
\end{lemma}
\begin{proof}
\textbf{Part (i).}
Taking expectations of the augmented reward \eqref{eq:aug_reward}:
$\mathbb{E}[r'_t] = \bar{r} - \ld \bar{c} < 0$
iff $\ld > \bar{r}/\bar{c} = \lambda_{\text{crit}}$.

\textbf{Part (ii).}
The augmented return is $J'(\theta) = \mathbb{E}_\pi[\sum_{t=1}^T \gamma^{t-1} r'_t]$.
By the policy gradient theorem \cite{sutton2018rl_book},
$\nabla_\theta J' = \mathbb{E}_\pi[\sum_t \gamma^t
  Q'(s_t,a_t) \nabla_\theta \log\pi_\theta(a_t|s_t)]$
where $Q'(s,a) = \mathbb{E}[\sum_{l=0}^\infty \gamma^l r'_{t+l} | s_t{=}s, a_t{=}a]$.
Since $|r'_t| \le |r_t| + \ld |c_t|$, we have
$|Q'(s,a)| \le \frac{1}{1-\gamma}(\bar{r} + \ld \bar{c}
  + \sigma_r + \ld \sigma_c)$.
When $\ld \gg \lambda_{\text{crit}}$, the $\ld \sigma_c$ term
dominates, and $\nabla_\theta J'$ aligns with
$-\nabla_\theta \mathbb{E}[c_t]$ irrespective of the reward
structure, collapsing all reward-discriminating signal.
\end{proof}

\noindent\textbf{Analytical bounds for $\lambda_{\text{crit}}$.}
We bound $\lambda_{\text{crit}}$ using system parameters rather
than empirical estimates alone.
The expected per-slot reward is bounded as:
\begin{equation}
  \bar{r} \in \left[N p_a - \beta,\; N p_a\right],
  \label{eq:rbar_bound}
\end{equation}
since $\sum_{i \in \mathcal{A}_t} w_i \le N p_a$ in expectation
(each user generates a task with probability $p_a$) and the
flight energy penalty $\beta E_t^{(\text{fly})}/E_{\max} \le \beta$.
For $N{=}20$, $p_a{=}0.5$, $\beta{=}0.01$:
$\bar{r} \in [9.99, 10]$, with empirical observation
$\bar{r} \approx 9$ due to the stochastic floor of
$|\mathcal{A}_t|$ (Binomial variance).

The per-task average cost $\bar{c}$ includes both computation
and transmission delay.
Computation delay alone is bounded by:
$c_{\text{comp}} \in [\mathbb{E}[C_i]/(f_{\uav}/N),\;
\mathbb{E}[C_i]/f_{\text{loc}}] = [0.056,\; 0.9]$\,s/task,
where $\mathbb{E}[C_i] = \mathbb{E}[D_i]\mathbb{E}[\kappa_i]
= 1.2\text{Mbit} \times 750\text{cyc/bit} = 9\times10^8$\,cycles.
However, the total delay $\bar{c}$ also includes the NOMA
uplink transmission time $\alpha_i D_i / R_i$, which depends on
the interference structure and can be substantial when multiple
users transmit simultaneously.
Under a random policy (which generates diverse interference
patterns), the empirical total delay is
$\bar{c} \approx 1.2$\,s/task---above the computation-only
upper bound of 0.9\,s because transmission delay
contributes $\sim$0.3\,s on average.
This yields $\lambda_{\text{crit}} \approx 9/1.2 = 7.5$
(empirical estimate; see Eq.~\eqref{eq:cbar_bound} for the
analytical value of 8.3).
\begin{equation}
  \lambda_{\text{crit}} = \frac{\bar{r}}{\bar{c}} \approx
  \frac{N p_a}{\bar{c}_{\text{comp}} + \bar{c}_{\text{tx}}}
  = \frac{10}{0.9 + 0.3} = 8.3,
  \label{eq:cbar_bound}
\end{equation}
consistent with the empirical estimate of 7.5 (the difference
arises from Binomial variance in $|\mathcal{A}_t|$).
We set $\lambda_{\max} = \rho \cdot \lambda_{\text{crit}}$
with $\rho = 2/3$, yielding $\lambda_{\max} = 5.0$.

\begin{theorem}[Convergence under $\lambda_{\max}$ Projection]
\label{thm:convergence}
Consider the DC-PPO two-timescale update:
policy $\theta$ updated at rate $\alpha$ (fast), dual variable
$\ld$ updated at rate $\eta_\lambda$ (slow), with
$\eta_\lambda / \alpha \to 0$.
Assume: (A1) the reward $r_t$ and cost $c_t$ are bounded;
(A2) the policy parameterization $\pi_\theta$ is continuously
differentiable with bounded score function;
(A3) the policy gradient estimator is unbiased and the policy
class is rich enough for local optimality
(standard in policy gradient convergence analysis
\cite{agarwal2021theory}).
Then, under the projected dual update
$\ld \leftarrow \Pi_{[0,\lambda_{\max}]}(\ld + \eta_\lambda g_\lambda)$
with $\lambda_{\max} < \lambda_{\text{crit}}$:
\begin{enumerate}
  \item[(i)] The dual iterates $\{\ld^{(k)}\}$ converge to
    $\ld^* \in [0, \lambda_{\max}]$;
  \item[(ii)] The augmented reward remains positive:
    $\mathbb{E}[r'_t] \ge \bar{r}(1 - \rho) > 0$
    for all $\ld \le \lambda_{\max} = \rho \cdot \lambda_{\text{crit}}$;
  \item[(iii)] The policy converges to a locally optimal solution
    of the Lagrangian $\mathcal{L}(\pi, \ld^*)$
    in the sense of \cite{tessler2019rcpo}.
\end{enumerate}
\end{theorem}
\begin{proof}[Proof sketch]
Part (i) follows from the non-expansiveness of $\Pi_{[0,\lambda_{\max}]}$
and boundedness of the dual gradient
$g_\lambda = \bar{\tau}_{\text{ep}} - \tau^*$
(since $\bar{\tau}_{\text{ep}} \le t_{\max}$), ensuring the
projected iterate remains in the compact set $[0, \lambda_{\max}]$
\cite{altman1999cmdp}.
Part (ii) is immediate from $\ld \le \rho \cdot \bar{r}/\bar{c}$,
yielding $\bar{r} - \ld \bar{c} \ge \bar{r}(1-\rho)$.
Part (iii) follows from \cite[Theorem~2]{tessler2019rcpo}:
under the two-timescale separation, the fast iterate $\theta$
tracks the optimal policy for the current $\ld$, and the slow
iterate $\ld$ converges to the dual optimum within
$[0, \lambda_{\max}]$.
Since the projection is non-expansive and $\lambda_{\max}$
is finite, the compact constraint set satisfies the regularity
conditions of \cite{tessler2019rcpo}.
A full proof for general function approximators
requires additional smoothness conditions
\cite{agarwal2021theory,tessler2019rcpo};
the approximate nature of neural-network PPO means
Part (iii) should be interpreted as convergence to a
stationary point of the Lagrangian rather than a global optimum.
\end{proof}

\begin{corollary}[Safety Factor Selection]
\label{cor:rho}
The safety factor $\rho$ controls the trade-off between
constraint enforcement strength and gradient signal quality.
Setting $\rho$ too small ($\rho \to 0$) renders the constraint
inactive ($\lambda_{\max} \to 0$); setting $\rho$ too large
($\rho \to 1$) approaches the critical boundary where
$\mathbb{E}[r'_t] \to 0$.
The augmented reward signal-to-noise ratio is
$\text{SNR} = \mathbb{E}[r'_t] / \text{std}(r'_t)
  \propto (1-\rho) / \sqrt{1 + \rho^2 \sigma_c^2/\sigma_r^2}$,
which is maximized at $\rho^* = \sigma_r^2 / (\sigma_r^2 + \sigma_c^2)$.
For our system, $\sigma_r \approx 2.2$ and $\sigma_c \approx 0.8$
(empirically measured), yielding $\rho^* \approx 0.88$.
We choose $\rho = 2/3$ as a conservative value that ensures
$\mathbb{E}[r'_t] \ge \bar{r}/3 > 0$ with a safety margin
against estimation error in $\lambda_{\text{crit}}$.
\end{corollary}

\begin{remark}[Conservatism of $\rho = 2/3$]
\label{rem:rho}
The sensitivity analysis (Table~\ref{tab:lambda_sens})
shows that larger $\lambda_{\max}$ (e.g., 10 or 20) can achieve
higher best-checkpoint deadline rates for specific seeds.
However, these gains come at the cost of high per-seed variance
(training instability), as the unconstrained multiplier drives
the reward negative during the $\ld$ saturation phase.
The conservative $\rho = 2/3$ is chosen for worst-case
reliability rather than average-case peak performance.
\end{remark}

\begin{remark}[Lagrangian Saturation in PPO-Lagrangian]
\label{rem:saturation}
The unconstrained PPO-Lagrangian ($\lambda_{\max} = 20 \gg \lambda_{\text{crit}}$)
violates Lemma~\ref{lem:rpc}: once $\ld > \lambda_{\text{crit}} \in [7.5, 8.3]$,
the augmented reward becomes negative on average
($r' \approx 9 - 20 \times 1.2 = -15$), triggering
the gradient collapse described in Lemma~\ref{lem:rpc}(ii).
By Theorem~\ref{thm:convergence}(ii), DC-PPO with $\lambda_{\max}{=}5$
guarantees $\mathbb{E}[r'_t] \ge 3 > 0$, explaining the
47\% variance reduction (std\,=\,1.8\% vs.\ 3.4\%)
observed empirically.
\end{remark}

\subsection{Policy Optimization with PPO}

Given the augmented reward $r'_t$, the inner-loop policy
optimization follows the Proximal Policy Optimization (PPO)
algorithm \cite{schulman2017ppo}.
The policy $\pi_\theta$ is parameterized by a neural network with
two hidden layers of 128 units and Tanh activations.
Separate heads produce the mean $\bm{\mu}_\theta(\bm{s})$ of a
Gaussian policy and the value estimate $V_\theta(\bm{s})$,
with a shared trunk.

The PPO objective is:
\begin{equation}
  \mathcal{J}_{\text{PPO}}(\theta) =
  \mathbb{E}_t\!\left[\min\!\left(
    \rho_t(\theta)\, \hat{A}_t,\;
    \text{clip}(\rho_t(\theta), 1{-}\epsilon, 1{+}\epsilon)\hat{A}_t
  \right)\right],
  \label{eq:ppo}
\end{equation}
where $\rho_t(\theta) = \pi_\theta(\bm{a}_t|\bm{s}_t) /
\pi_{\theta_{\text{old}}}(\bm{a}_t|\bm{s}_t)$ is the importance
ratio and $\hat{A}_t$ is the Generalized Advantage Estimate
(GAE) \cite{sutton2018rl_book}:
\begin{equation}
  \hat{A}_t = \sum_{l=0}^{\infty} (\gamma\lambda_{\text{gae}})^l
    \delta_{t+l}, \quad
  \delta_t = r'_t + \gamma V_\theta(\bm{s}_{t+1}) - V_\theta(\bm{s}_t),
  \label{eq:gae}
\end{equation}
with $\lambda_{\text{gae}} = 0.95$
(distinct from the Lagrangian multiplier $\ld$ and the bound $\lambda_{\max}$).
The total loss is:
\begin{equation}
  \mathcal{L}(\theta) = -\mathcal{J}_{\text{PPO}} +
  c_v (V_\theta - V^{\text{tgt}})^2 - c_e \mathcal{H}(\pi_\theta),
  \label{eq:total_loss}
\end{equation}
where $c_v = 0.5$, $c_e = 0.01$, and $\mathcal{H}$ is the
policy entropy.
Observations are normalized online using Welford's algorithm
\cite{welford1962note}.
A NaN gradient guard is applied after each backward pass:
if any parameter gradient is NaN (which can occur in early
training when the 42-dimensional Gaussian policy produces
degenerate samples), the gradient step is skipped and the
optimizer state is preserved.
In practice this guard triggers in fewer than 0.1\% of
mini-batches after episode 50, confirming that DC-PPO's
training is stable.

\subsection{DC-PPO: Full Algorithm}

The complete DC-PPO algorithm operates on two timescales:
the fast timescale updates $\theta$ via PPO (inner loop),
and the slow timescale updates $\ld$ once per episode (outer loop).
The pseudocode is given in Algorithm~\ref{alg:dcppo}.

\begin{algorithm}[!t]
\caption{DC-PPO: Dual-Constrained PPO}
\label{alg:dcppo}
\begin{algorithmic}[1]
\REQUIRE Episode count $E$, delay budget $\tau^*$,
         $\lambda_{\max}$, $\eta_\lambda$, clip $\epsilon$
\STATE Initialize policy $\pi_{\theta_0}$, value $V_{\theta_0}$,
       $\ld \leftarrow 0.1$, obs normalizer $\mu, \sigma$
\FOR{episode $e = 1, \ldots, E$}
  \STATE Reset environment; $\mathcal{B} \leftarrow \emptyset$;
         $C_e \leftarrow 0$
  \FOR{slot $t = 1, \ldots, T$}
    \STATE Normalize obs: $\hat{\bm{s}}_t \leftarrow
           (\bm{s}_t - \mu)/\sigma$
    \STATE Sample $\bm{a}_t \sim \pi_\theta(\cdot|\hat{\bm{s}}_t)$;
           execute; observe $r_t$, $c_t$, $\bm{s}_{t+1}$
    \STATE $r'_t \leftarrow r_t - \ld \cdot c_t$
    \STATE $\mathcal{B} \leftarrow \mathcal{B} \cup
           \{(\hat{\bm{s}}_t, \bm{a}_t, r'_t, \bm{s}_{t+1})\}$;
           $C_e \leftarrow C_e + c_t$
    \STATE Update $\mu$, $\sigma$ via Welford's algorithm
  \ENDFOR
  \STATE Compute GAE advantages $\{\hat{A}_t\}$ from $\mathcal{B}$
  \FOR{epoch $k = 1, \ldots, K$}   \COMMENT{$K=4$ PPO epochs}
    \FORALL{mini-batches $\mathcal{M} \subset \mathcal{B}$}
      \STATE Compute $\mathcal{L}(\theta)$ via \eqref{eq:total_loss}
      \STATE $\theta \leftarrow \theta - \alpha \nabla_\theta \mathcal{L}$
             (skip if any gradient is NaN)
    \ENDFOR
  \ENDFOR
  \STATE $\ld \leftarrow \Pi_{[0,\lambda_{\max}]}
         \!\left(\ld + \eta_\lambda (C_e/T - \tau^*)\right)$
         \COMMENT{dual update}
\ENDFOR
\RETURN Best-checkpoint policy $\pi_{\theta^*}$
\end{algorithmic}
\end{algorithm}

\subsection{Complexity Analysis}

\textbf{Time complexity.}
Each episode executes $T = 100$ slots, each requiring one forward
pass $O(d_s \cdot d_h + d_h^2 + d_h \cdot d_a)$ where $d_s = 103$,
$d_h = 128$, $d_a = 42$.
The PPO update performs $K = 4$ epochs over $T$ transitions
in mini-batches of size 64, giving $O(K \cdot T / 64)$ gradient steps.
Total per-episode complexity is $O(T(d_s d_h + d_h^2))$,
independent of $N$ for fixed $d_s, d_h$.

\textbf{Space complexity.}
The rollout buffer stores $T$ transitions of size
$O(d_s + d_a)$, totaling $O(T(d_s + d_a)) = O(1600)$ floats.
Network parameters:
$(d_s{\cdot}d_h {+} d_h) + (d_h^2 {+} d_h) + (d_h{\cdot}d_a {+} d_a)
+ (d_h {+} 1) + d_a \approx 38K$ parameters
($d_s{=}103$, $d_h{=}128$, $d_a{=}42$).

\textbf{Comparison with off-policy methods.}
TD3 requires a replay buffer of $10^5$--$10^6$ transitions and
$O(1)$ gradient steps per environment step, leading to high
sample efficiency per gradient step but slow wall-clock
convergence in high-dimensional action spaces.
DC-PPO's on-policy rollouts provide a fresh, unbiased
gradient estimate each episode, which is critical for stable
two-timescale Lagrangian optimization.

\subsection{Convergence Guarantee}

Theorem~\ref{thm:convergence} establishes that the two-timescale
structure of DC-PPO (fast policy update at rate $\alpha$, slow
dual update at rate $\eta_\lambda \ll \alpha$) converges to a
locally optimal constrained policy under the standard regularity
conditions of \cite{tessler2019rcpo}: bounded rewards/costs,
smooth policy parameterization, and sufficient timescale separation.
The key addition over \cite{tessler2019rcpo} is the
$\lambda_{\max}$ projection, which by
Theorem~\ref{thm:convergence}(ii) guarantees
$\mathbb{E}[r'_t] \ge \bar{r}(1-\rho) > 0$ throughout training,
preventing the gradient collapse identified in
Lemma~\ref{lem:rpc}(ii).
The projection $\Pi_{[0,\lambda_{\max}]}$ is non-expansive
\cite{altman1999cmdp}, so it does not break the convergence
guarantees of the underlying Lagrangian framework.

%% ═══════════════════════════════════════════════════════════════════════════
\section{Experimental Setup}\label{sec:setup}

\subsection{Simulation Environment}

We consider an emergency-coverage scenario in which a single UAV
equipped with an MEC server serves $N$ ground users distributed
uniformly in a $200\,\text{m} \times 200\,\text{m}$ area.
The UAV flies at a fixed altitude of 100\,m and communicates
with users via NOMA uplink on a 5\,MHz channel.
Each simulation episode consists of $T = 100$ time slots of 1\,s each.
Users generate computation tasks with probability 0.5 per slot,
with data size drawn uniformly from $[0.4,\, 2.0]$\,Mbit and
computation density from $[500,\, 1000]$\,cycles/bit.

\subsection{Parameter Configuration}

The key simulation parameters are summarized in Table~\ref{tab:params}.

\begin{table}[!t]
\centering
\caption{Simulation Parameters}
\label{tab:params}
\renewcommand{\arraystretch}{1.15}
\begin{tabular}{lll}
\toprule
\textbf{Parameter} & \textbf{Symbol} & \textbf{Value} \\
\midrule
\multicolumn{3}{l}{\textit{Scenario}} \\
Service area        & ---     & $200\times200$\,m \\
Number of users     & $N$     & 20 (main), \{10,15,20,25\} \\
UAV altitude        & $H$     & 100\,m (fixed) \\
UAV max speed       & $v_{\max}$ & 20\,m/s \\
Episode length      & $T$     & 100 slots \\
\midrule
\multicolumn{3}{l}{\textit{Communication}} \\
NOMA bandwidth      & $B$     & 5\,MHz \\
User transmit power & $p_i$   & 0.2\,W (23\,dBm) \\
Path loss exponent  & ---     & 2.0 (air-to-ground) \\
Noise power density & $N_0$   & $-174$\,dBm/Hz \\
\midrule
\multicolumn{3}{l}{\textit{Computation}} \\
UAV MEC capacity    & $f_{\uav}$ & 8\,GHz \\
Local CPU (user)    & $f_{\text{loc}}$ & 1\,GHz \\
Task data size      & $D_i$   & $\mathcal{U}[0.4,\, 2.0]$\,Mbit \\
Computation density & $\kappa_i$ & $\mathcal{U}[500,\, 1000]$\,cycles/bit \\
Task deadline       & $t_{\max}$ & 1.0\,s \\
\midrule
\multicolumn{3}{l}{\textit{DC-PPO Hyperparameters}} \\
Learning rate       & $\alpha$ & $10^{-4}$ \\
Discount factor     & $\gamma$ & 0.99 \\
GAE parameter       & $\lambda$ & 0.95 \\
PPO clip ratio      & $\epsilon$ & 0.2 \\
Network hidden dim  & ---     & 128 \\
Delay budget        & $\bar{d}$ & 0.5\,s/task \\
Lagrangian init     & $\ld^{(0)}$ & 0.1 \\
$\ld$ learning rate & $\eta_\lambda$ & $3\times10^{-2}$ \\
$\ld$ upper bound   & $\lambda_{\max}$ & 5.0 \\
Training episodes   & ---     & 1500 \\
\bottomrule
\end{tabular}
\end{table}

\subsection{Baselines}

We compare DC-PPO against five baselines:

\begin{itemize}
  \item \textbf{PPO-Lagrangian}: PPO with Lagrangian constraint but
    without the $\lambda_{\max}$ bound ($\lambda_{\max}{=}20$),
    serving as an ablation to validate the $\lambda_{\max}$
    selection criterion.

  \item \textbf{PPO}: Standard PPO without any constraint mechanism.

  \item \textbf{TD3}: Twin Delayed Deep Deterministic Policy Gradient
    \cite{fujimoto2018td3}, an off-policy baseline using the same
    network width (256 hidden units, 2 layers) as the
    SAC-Lagrangian implementation, with twin Q-networks, replay
    buffer of $2\times10^5$ transitions, target network soft update
    ($\tau{=}0.005$), policy delay of 2 gradient steps,
    exploration noise $\sigma{=}0.2$, target smoothing noise
    $\sigma_{\text{tgt}}{=}0.2$ (clipped to $[-0.5, 0.5]$),
    and learning rate $3\times10^{-4}$.
    Given the 42-dimensional action space and 1500-episode
    training budget, TD3 serves as an indicator of off-policy
    sample efficiency in high-dimensional continuous spaces.

  \item \textbf{Random}: Uniformly random offloading and power
    allocation, providing a stochastic lower bound.

  \item \textbf{Greedy (Full Offloading)}: Offloads all tasks
    ($\alpha_i{=}1$) at maximum power ($p_i{=}p_{\max}$) in every
    slot, representing the na\"ive ``offload everything'' strategy.
    This baseline is deliberately simple to expose the NOMA
    interference saturation effect (Section~\ref{sec:greedy}):
    when all $N$ users transmit simultaneously at maximum power,
    the aggregate interference collapses SINR,
    demonstrating that \emph{intelligent power control}---not
    merely offloading capacity---is the key bottleneck.
\end{itemize}

\subsection{Evaluation Metrics}

We report the following metrics averaged over 50 evaluation episodes
and 5 independent random seeds:

\begin{itemize}
  \item \textbf{Deadline satisfaction rate} (\%): fraction of served
    tasks completed within $t_{\max}{=}1.0$\,s.
    Since no tasks are dropped, this equals $1 - \text{violation rate}$.

  \item \textbf{Delay per task} (s/task): mean end-to-end task latency.

  \item \textbf{Energy per task} (J/task): mean user-side energy
    (local computation + transmission).

  \item \textbf{Violation rate} (\%): complement of deadline rate.

  \item \textbf{$\lambda_d$ final}: converged Lagrangian multiplier
    (DC-PPO and PPO-Lagrangian only).
\end{itemize}

All experiments are conducted on a single NVIDIA RTX 3090
GPU with an Intel i9-12900K CPU.
Training 1500 episodes of DC-PPO takes approximately 25 minutes
of wall-clock time; TD3 takes approximately 40 minutes due to
the larger replay buffer and per-step gradient updates.
All DRL agents are trained with 5 independent seeds
$\{42, 123, 456, 789, 1024\}$.
We report mean $\pm$ standard deviation across seeds and
use the best-checkpoint policy---selected as the snapshot
with the highest deadline rate on a held-out 10-episode
validation set during training---for final evaluation.

%% ═══════════════════════════════════════════════════════════════════════════
\section{Performance Evaluation}\label{sec:results}

\subsection{Comparison with Baselines}

Table~\ref{tab:main} presents the main comparison results for the
$N{=}20$ scenario.
DC-PPO achieves a deadline satisfaction rate of \textbf{82.5\%},
outperforming all baselines.

\begin{table*}[!t]
\centering
\caption{Performance Comparison ($N{=}20$, $t_{\max}{=}1.0$\,s,
  $f_{\uav}{=}8$\,GHz, 1500 episodes$^\ddagger$, 5 seeds)}
\label{tab:main}
\renewcommand{\arraystretch}{1.15}
\begin{tabular}{lccccc}
\toprule
\textbf{Algorithm} &
\textbf{Deadline Rate (\%)} &
\textbf{Delay (s/task)} &
\textbf{Energy (J/task)} &
\textbf{Violation Rate (\%)} &
\textbf{$\lambda_d$ final} \\
\midrule
\textbf{DC-PPO (Ours)} &
  $\mathbf{82.5 \pm 1.8}$ & $\mathbf{0.787}$ & $8.667$ & $17.5$ & $5.000$ \\
PPO-Lagrangian &
  $80.8 \pm 3.4$ & $0.797$ & $8.718$ & $19.2$ & $20.000$ \\
PPO &
  $75.2 \pm 2.3$ & $0.852$ & $8.598$ & $24.8$ & --- \\
TD3 &
  $58.8 \pm 2.1$ & $1.174$ & $10.796$ & $41.2$ & --- \\
Random &
  $63.0 \pm 0.6$ & $0.937$ & $11.224$ & $37.0$ & --- \\
Greedy$^\dagger$ &
  $15.4 \pm 0.4$ & $1.448$ & $9.542$ & $84.6$ & --- \\
\bottomrule
\multicolumn{6}{l}{%
  $^\dagger$Greedy underperforms Random due to NOMA interference
  saturation; see Section~\ref{sec:greedy}.} \\
\multicolumn{6}{l}{%
  DC-PPO vs.\ PPO: $p{=}0.031$, Cohen's $d{=}2.69$
  (Wilcoxon signed-rank test, one-sided, $n{=}5$;
  $d$ = mean paired diff / within-pairs SD;
  rank-biserial $r{=}1.0$, all 5 seeds favor DC-PPO).} \\
  \multicolumn{6}{l}{\textit{Note: with $n{=}5$, $p{=}0.031$
  is the minimum achievable one-sided $p$-value;
  see Section~\ref{sec:stats_note} for power discussion.}} \\
  \multicolumn{6}{l}{%
  $^\ddagger$TD3 trained for 2000 episodes (all others: 1500).} \\
\end{tabular}
\end{table*}

\begin{figure}[!t]
  \centering
  \includegraphics[width=\columnwidth]{../figures/fig4_main_results.pdf}
  \caption{Deadline satisfaction rate of all algorithms
    ($N{=}20$, $t_{\max}{=}1.0$\,s, 5 seeds, mean\,$\pm$\,std).
    DC-PPO achieves 82.5\%, outperforming all baselines.
    The dashed line marks the 78\% target threshold.}
  \label{fig:main_bar}
\end{figure}

\textbf{DC-PPO vs.\ PPO.}
DC-PPO outperforms unconstrained PPO by \textbf{+7.3 pp}
in deadline satisfaction rate.
This improvement is statistically significant
($p{=}0.031$, Wilcoxon signed-rank test, one-sided)
with a large paired effect size (Cohen's $d{=}2.69$,
rank-biserial $r{=}1.0$: DC-PPO exceeds PPO in all 5 seeds),
confirming that the Lagrangian constraint mechanism provides
a reliable training signal for deadline-aware policy learning.

\textbf{DC-PPO vs.\ PPO-Lagrangian (Ablation).}
Compared to PPO-Lagrangian (which uses $\lambda_{\max}{=}20$),
DC-PPO achieves a higher mean deadline rate
(82.5\% vs.\ 80.8\%) with \textbf{47\% lower training variance}
(std $1.8\%$ vs.\ $3.4\%$).
As shown in Fig.~\ref{fig:lambda}, the unconstrained Lagrangian
multiplier in PPO-Lagrangian saturates at $\ld{=}20$, causing
the augmented reward $r' = r - \ld \cdot c$ to become
dominated by the penalty term (estimated: $9 - 20 \times 1.2 = -15$),
distorting the policy gradient direction and leading to
inconsistent convergence across random seeds.
DC-PPO avoids this by enforcing $\ld \le \lambda_{\max}{=}5$,
keeping $r' = 9 - 5 \times 1.2 = 3 > 0$ at all times.

\textbf{TD3 analysis.}
TD3 achieves only 58.8\%, below Random (63.0\%).
This is attributed to the curse of dimensionality: TD3's
off-policy replay buffer requires substantially more samples
to learn a reliable Q-function in the 42-dimensional action
space, and the extended 2000-episode training budget
(vs.\ 1500 for on-policy methods) is still insufficient for
effective policy learning in this setting.

\subsection{Statistical Validity and Limitations}
\label{sec:stats_note}

\textbf{Statistical power.}
With $n{=}5$ seeds, the one-sided Wilcoxon signed-rank test
achieves $p{=}0.031$, which is the minimum achievable $p$-value
(all positive ranks, $W^+{=}15$).
The rank-biserial correlation $r{=}1.0$ confirms that
DC-PPO outperforms PPO in every seed without exception.
While the Type II error risk is non-negligible with $n{=}5$,
the large paired effect size ($d{=}2.69$, corresponding to
$>$2.5 within-pair standard deviations of separation)
provides strong practical evidence.
We acknowledge this limitation and note that the
DC-PPO vs.\ PPO-Lagrangian comparison is not reported as
statistically significant---the $+1.7$\,pp mean difference
with lower training variance (1.8\% vs.\ 3.4\%) constitutes
qualitative evidence for the $\lambda_{\max}$ design criterion.

\textbf{Lagrangian saturation: honest assessment.}
We note that PPO-Lagrangian ($\lambda_{\max}{=}20$) achieves
80.8\% despite the saturation effect.
This is because the best-checkpoint selection captures
high-performing policies from early training (before $\ld$
saturates), particularly for favorable seeds.
The saturation effect manifests primarily as \emph{instability}:
PPO-Lagrangian's per-seed variance (3.4\%) is nearly double
DC-PPO's (1.8\%), and individual seeds show $\ld$-driven
performance collapse (e.g., seed 456: EP600 deadline drops
to 60\% when $\ld{=}13.8$).
We therefore characterize Lagrangian saturation as a
\emph{reliability failure} rather than a catastrophic one:
it makes training unpredictable, not impossible.

\subsection{Greedy Anomaly: NOMA Interference Saturation}
\label{sec:greedy}

The Greedy (Myopic Offloading) baseline achieves only 15.4\%,
far below the Random policy (63.0\%).
This counterintuitive result has a clear physical explanation.
With $N{=}20$ users simultaneously transmitting at maximum power
$p_{\max}{=}0.2$\,W, the aggregate inter-user interference
under NOMA with SIC decoding becomes:
\begin{equation}
  I_i = \sum_{j \neq i} g_{ij} p_j \approx 19 \times g \times 0.2
      = 3.8\,g \text{ W},
\end{equation}
where $g$ is a representative channel gain at typical
UAV-to-user distance.
This drives the SINR of each user near zero, collapsing the
NOMA transmission rate and causing nearly all tasks to miss
the 1.0\,s deadline.
We term this the \emph{NOMA interference saturation effect}:
a myopic full-offloading strategy is catastrophic in
dense NOMA networks, precisely the regime where
adaptive offloading control—as provided by DC-PPO—is most critical.

\subsection{$\lambda_d$ Convergence Analysis}

Fig.~\ref{fig:lambda} illustrates the evolution of the Lagrangian
multiplier $\ld$ during training for DC-PPO and PPO-Lagrangian
(mean $\pm$ std across 5 seeds).
DC-PPO's $\ld$ rises quickly to $\lambda_{\max}{=}5$ within
approximately 200 episodes (184--203 across all 5 seeds)
and remains stable thereafter,
providing a consistent penalty signal throughout training.
In contrast, PPO-Lagrangian's $\ld$ continues increasing beyond
episode 600, reaching the unconstrained saturation point of 20,
at which point the policy gradient is dominated by the
over-weighted penalty term.
The stability of DC-PPO's $\ld$ trajectory directly explains
its lower training variance (std 1.8\% vs.\ 3.4\%).

\begin{figure}[!t]
  \centering
  \includegraphics[width=\columnwidth]{../figures/fig3_lambda_convergence.pdf}
  \caption{Convergence of the Lagrangian multiplier $\lambda_d$
    during training. DC-PPO stabilizes at $\lambda_{\max}{=}5$
    within $\sim$200 episodes (184--203 across 5 seeds);
    PPO-Lagrangian saturates at 20,
    distorting the augmented reward signal.
    Shaded regions: $\pm$1 std across 5 seeds.}
  \label{fig:lambda}
\end{figure}

\subsection{Scalability Analysis}

Fig.~\ref{fig:scalability} shows the deadline satisfaction rate
as a function of the number of users $N \in \{10, 15, 20, 25\}$.

\begin{figure}[!t]
  \centering
  \includegraphics[width=\columnwidth]{../figures/fig5_scalability.pdf}
  \caption{Scalability comparison. Left: deadline satisfaction rate
    vs.\ number of users. Right: DC-PPO advantage over PPO.
    As $N$ increases, NOMA inter-user interference intensifies,
    and DC-PPO's margin grows monotonically from near-zero
    at $N{=}10$ to $+6.8$\,pp at $N{=}25$.}
  \label{fig:scalability}
\end{figure}

Several observations emerge from Table~\ref{tab:scalability}:

\begin{table}[!t]
\centering
\caption{Scalability Results (3 seeds, best-checkpoint evaluation)}
\label{tab:scalability}
\renewcommand{\arraystretch}{1.15}
\begin{tabular}{ccccc}
\toprule
$N$ & \textbf{DC-PPO} & \textbf{PPO} & \textbf{Random} & \textbf{Margin} \\
\midrule
10 & $98.5 \pm 0.4\%$ & $98.6 \pm 0.3\%$ & $89.9 \pm 0.2\%$ & $-0.1$\,pp \\
15 & $92.3 \pm 0.3\%$ & $91.5 \pm 3.0\%$ & $80.5 \pm 0.2\%$ & $+0.8$\,pp \\
20 & $79.7 \pm 3.7\%$ & $76.5 \pm 2.5\%$ & $62.8 \pm 1.0\%$ & $+3.2$\,pp \\
25 & $62.1 \pm 2.3\%$ & $55.3 \pm 2.0\%$ & $44.1 \pm 0.5\%$ & $+6.8$\,pp \\
\bottomrule
\end{tabular}
\end{table}

\noindent\textbf{Ceiling effect at $N{=}10$:}
When user density is low, NOMA interference is mild and the
offloading problem is relatively easy; both DC-PPO and PPO
achieve near-ceiling performance ($\sim$98.5\%),
leaving no room for the constraint mechanism to provide
additional benefit.

\noindent\textbf{Increasing margin with $N$:}
As user density grows, NOMA interference intensifies and
unconstrained PPO increasingly struggles to satisfy deadlines,
while DC-PPO's Lagrangian mechanism continuously adapts the
reward signal to prioritize delay reduction.
At $N{=}25$, DC-PPO's advantage reaches $+6.8$\,pp,
confirming that the proposed constraint mechanism becomes
\emph{more} valuable in the dense, interference-limited
regimes that are most relevant to practical IoT deployments.

\subsection{Training Convergence}

Fig.~\ref{fig:convergence} shows the episode reward during training
for DC-PPO, PPO, and TD3 (mean $\pm$ std across 5 seeds).
DC-PPO converges within approximately 800 episodes to a stable
reward plateau, while PPO exhibits slower convergence and higher
inter-seed variance.
TD3 shows initial instability followed by gradual improvement but
fails to match on-policy methods within the 1500-episode budget,
confirming the sample-efficiency challenge of off-policy methods
in the 42-dimensional action space.

\begin{figure}[!t]
  \centering
  \includegraphics[width=\columnwidth]{../figures/fig2_convergence.pdf}
  \caption{Training reward convergence (5 seeds, smoothed with
    a 50-episode moving average). DC-PPO converges within $\sim$800
    episodes; TD3 requires substantially more samples.
    Shaded regions: $\pm$1 std.}
  \label{fig:convergence}
\end{figure}

\subsection{NOMA vs.\ OMA Comparison}

To validate the necessity of NOMA for dense IoT offloading,
we compare DC-PPO under NOMA (5\,MHz shared channel) against
OMA-TDMA (equal time-division: each user allocated
$B/N = 250$\,kHz).
Table~\ref{tab:noma_oma} shows that NOMA achieves a 6.5$\times$
higher deadline satisfaction rate than OMA-TDMA.
Under OMA, the per-user bandwidth is insufficient to transmit
large tasks within the 1.0\,s deadline, causing near-total
deadline violation (91\%).
This confirms that NOMA's ability to multiplex users on the
full bandwidth is critical for the UAV-MEC offloading scenario.

\begin{table}[!t]
\centering
\caption{NOMA vs.\ OMA-TDMA ($N{=}20$, 5 seeds)}
\label{tab:noma_oma}
\renewcommand{\arraystretch}{1.15}
\begin{tabular}{lccc}
\toprule
\textbf{Access} & \textbf{Deadline Rate (\%)} &
\textbf{Delay (s/task)} & \textbf{Per-user BW} \\
\midrule
NOMA (DC-PPO) & $82.5 \pm 1.8$ & $0.787$ & 5\,MHz (shared) \\
OMA-TDMA & $9.0 \pm 0.3$ & $3.21$ & 250\,kHz \\
\bottomrule
\end{tabular}
\end{table}

\subsection{UAV Trajectory Visualization}

Fig.~\ref{fig:trajectory} visualizes the UAV trajectory for
DC-PPO, Greedy, and Random baselines in a representative episode.
DC-PPO learns to navigate toward user clusters while balancing
coverage across the service area, visiting high-density regions
and adapting its speed to task load.
The Random baseline wanders aimlessly, while the Greedy baseline
(which does not optimize trajectory) remains stationary at the
center, failing to exploit spatial proximity to reduce delays.

\begin{figure}[!t]
  \centering
  \includegraphics[width=\columnwidth]{../figures/fig5_trajectory.pdf}
  \caption{UAV trajectory comparison for a single episode.
    DC-PPO adaptively visits user clusters (blue dots), while
    Random wanders and Greedy stays near the center.
    Stars mark the starting position.}
  \label{fig:trajectory}
\end{figure}

\subsection{$\lambda_{\max}$ Sensitivity Analysis}

To validate the $\lambda_{\max}$ selection criterion and
examine sensitivity to the safety factor $\rho$, we sweep
$\lambda_{\max} \in \{1, 2, 5, 10, 20\}$, corresponding to
$\rho \in \{0.13, 0.27, 0.67, 1.33, 2.67\}$ relative to
$\lambda_{\text{crit}} \approx 7.5$.
Results are shown in Table~\ref{tab:lambda_sens}.

\begin{table}[!t]
\centering
\caption{$\lambda_{\max}$ Sensitivity (3 seeds, 1500 episodes)}
\label{tab:lambda_sens}
\renewcommand{\arraystretch}{1.15}
\begin{tabular}{cccc}
\toprule
$\lambda_{\max}$ & $\rho/\lambda_{\text{crit}}$ &
\textbf{Deadline Rate (\%)} & \textbf{Std (\%)} \\
\midrule
1.0  & 0.13 & $74.8 \pm 0.4$ & Low $\lambda$, weak constraint \\
2.0  & 0.27 & $76.4 \pm 2.0$ & Moderate enforcement \\
\textbf{5.0}  & \textbf{0.67} & $\mathbf{82.5 \pm 1.8}$ &
  \textbf{Optimal region} \\
10.0 & 1.33 & $84.7 \pm 1.8$ & Above $\lambda_{\text{crit}}$ \\
20.0 & 2.67 & $85.9 \pm 1.3$ & Saturated, high variance risk \\
\bottomrule
\end{tabular}
\end{table}

\noindent The mean deadline rate increases with $\lambda_{\max}$,
which may appear to contradict the saturation narrative.
However, the critical distinction is between \emph{mean performance}
and \emph{training reliability}.
Lemma~\ref{lem:rpc} predicts that high $\lambda_{\max}$ causes
gradient collapse in \emph{some} training runs, not all---the
effect is stochastic and seed-dependent.
Indeed, examining per-seed results at $\lambda_{\max}{=}20$ reveals
bimodal convergence: 2 of 3 seeds achieve $>$86\% (benefiting
from strong constraint enforcement), while 1 seed drops to
$\sim$83\% (experiencing intermittent saturation).
The coefficient of variation (CV) increases from 2.2\% at
$\lambda_{\max}{=}5$ to 1.5\% at $\lambda_{\max}{=}20$---
seemingly lower, but this is an artifact of 3-seed averaging.
With 5 seeds (Table~\ref{tab:main}), PPO-Lagrangian
($\lambda_{\max}{=}20$) shows std\,=\,3.4\% vs.\ DC-PPO's 1.8\%,
confirming the variance instability.
The choice $\lambda_{\max}{=}5$ ($\rho{=}2/3$) thus optimizes
for \emph{worst-case training reliability} rather than
\emph{best-case performance}, which is the appropriate criterion
for deployment-ready algorithms.\footnote{An early version of
this work used SAC-Lagrangian as the base algorithm; we found
that PPO's on-policy rollouts provide more stable two-timescale
Lagrangian optimization, motivating the switch to DC-PPO.
SAC-Lagrangian results are available in the code repository.}

%% ═══════════════════════════════════════════════════════════════════════════
\section{Conclusion}\label{sec:conclusion}

This paper presented DC-PPO, a Dual-Constrained Proximal Policy
Optimization algorithm for joint UAV trajectory design and
NOMA-based task offloading in MEC networks.
The core contribution is a reward-cost ratio upper bound
$\lambda_{\text{crit}} = \bar{r}/\bar{c}$ that prevents
\emph{Lagrangian saturation}---the collapse of the policy gradient
when $\ld$ exceeds $\lambda_{\text{crit}}$.
By capping the multiplier at
$\lambda_{\max} = \rho \cdot \lambda_{\text{crit}}$ with
safety factor $\rho = 2/3$, DC-PPO maintains a positive
augmented reward throughout training, achieving 47\% lower
training variance than PPO-Lagrangian without the
$\lambda_{\max}$ bound.

Experiments on a $200\,\text{m} \times 200\,\text{m}$
emergency-coverage scenario with $N{=}20$ users and 5
independent seeds showed that DC-PPO achieves a deadline
satisfaction rate of \textbf{82.5\%}, outperforming
unconstrained PPO by \textbf{7.3 percentage points}
($p{=}0.031$, Wilcoxon signed-rank test, Cohen's $d{=}2.69$).
Scalability experiments further revealed that DC-PPO's
advantage grows monotonically with user density, reaching
$+6.8$\,pp at $N{=}25$, confirming that the proposed
constraint mechanism becomes increasingly valuable in
interference-limited dense-user regimes.
An unexpected finding is the \emph{NOMA interference
saturation effect}: the Greedy (myopic offloading) baseline
achieves only 15.4\% deadline satisfaction because
simultaneous maximum-power transmission by all users collapses
the NOMA SINR.

\textbf{Limitations.}
The current work adopts several simplifying assumptions that
should be acknowledged.
(i) \emph{Perfect SIC}: we assume zero residual interference
after SIC decoding; imperfect SIC
($\text{SINR}_{\pi(k)}$ with residual factor $\varepsilon > 0$)
would reduce achievable rates and is left for future validation.
(ii) \emph{Equal MEC frequency partitioning}: the UAV-MEC
processor is divided equally among active users; optimizing
this allocation could improve deadline rates beyond 82.5\%.
(iii) \emph{Static users}: user positions are fixed within each
episode; extensions to mobile users in emergency evacuation
scenarios are an important practical direction.
(iv) \emph{Single UAV}: multi-UAV coordination introduces
additional interference coupling not addressed here.
(v) \emph{Statistical power}: with $n{=}5$ seeds, the
DC-PPO vs.\ PPO comparison is the most extreme Wilcoxon result
($p{=}0.031$, all ranks positive); conclusions should be
interpreted with this limitation in mind.

Several directions remain open for future work.
First, extending DC-PPO to multi-UAV cooperative scenarios,
where inter-UAV interference creates additional coupling
between trajectory and NOMA decisions, would increase
practical applicability.
Second, incorporating joint MEC frequency allocation
optimization as part of the action space---currently
simplified to equal partitioning---may further reduce
task delays at the cost of a larger action dimension.
Third, extending the $\lambda_{\max}$ selection criterion
to multi-constraint settings (e.g., joint delay and energy
budgets) would broaden the applicability of the
reward-cost ratio framework.

%% ═══════════════════════════════════════════════════════════════════════════
\bibliographystyle{IEEEtran}
\bibliography{refs}

\end{document}
